\documentclass[10pt,a4paper]{amsart}
\usepackage[margin=19mm]{geometry}
\usepackage{amsmath,amssymb,mathtools}
\usepackage[scr=rsfs,cal=euler]{mathalfa}
\usepackage{xfrac}
\usepackage[unicode,hidelinks,bookmarks=false]{hyperref}
\newcommand{\ie}{i.e.\ }
\newcommand{\cf}{cf.\ }
\newcommand{\resp}{respectively\ }
\newcommand{\Subsection}{\subsection}
\providecommand{\nobreakdash}{\nobreak\hbox{-}}
\setlength{\emergencystretch}{2em}
\multlinegap0pt
\usepackage{amscd,mathrsfs}
\usepackage{verbatim}
\usepackage{enumitem}
\usepackage{tikz-cd}
\setlength{\textheight}{9in}\setlength{\textwidth}{475pt}
\newcommand{\newsection}[1]{\setcounter{equation}{0} \section{#1}}
\setcounter{footnote}{1}
\newtheorem{theorem}{Theorem}[section]
\newtheorem{lemma}[theorem]{Lemma}
\newtheorem{corollary}[theorem]{Corollary}
\newtheorem{proposition}[theorem]{Proposition}
\newtheorem{definition}[theorem]{Definition}
\newtheorem{example}[theorem]{Example}
\newtheorem{exercise}[theorem]{Exercise}
\newtheorem{remark}[theorem]{Remark}
\newtheorem{question}[theorem]{Question}
\newtheorem{conjecture}[theorem]{Conjecture}
\newtheorem{problem}[theorem]{Problem}
\def \D{\mathbb{D}}
\def \N{\mathbb{N}}
\def \Z{\mathbb{Z}}
\def \F{\mathbb{F}}
\def \R{\mathbb{R}}
\def \C{\mathbb{C}}
\def \T{\mathbb{T}}
\newcommand{\wt}{\widetilde}
\newcommand{\clb}{\mathcal{B}}
\newcommand{\clg}{\mathcal{G}}
\newcommand{\cle}{\mathcal{E}}
\newcommand{\clh}{\mathcal{H}}
\newcommand{\clm}{\mathcal{M}}
\newcommand{\cln}{\mathcal{N}}
\newcommand{\clr}{\mathcal{R}}
\newcommand{\clw}{\mathcal{W}}
\newcommand{\clk}{\mathcal{K}}
\newcommand{\cls}{\mathcal{S}}
\newcommand{\clf}{\mathcal{F}}
\newcommand{\clq}{\mathcal{Q}}
\newcommand{\cld}{\mathcal{D}}
\newcommand{\clo}{\mathcal{O}}
\newcommand{\cli}{\mathcal{I}}
\newcommand{\ran}{\mathrm{ran}}
\newcommand{\dist}{\mathrm{dist}}
\newcommand{\codim}{\mathrm{codim}}
\newcommand{\essran}{\mathrm{ess\ ran}}
\newcommand{\essinf}{\mathrm{ess\ inf}}
\newcommand{\esssup}{\mathrm{ess\ sup}}
\newcommand{\vp}{\varphi}
\newcommand{\ind}{\operatorname{ind}}
\newcommand{\diag}{\mathrm{diag}}
\newcommand{\ip}[2]{\left\langle #1,#2\right\rangle}
\newcommand{\norm}[1]{\left\|#1\right\|}
\newcommand{\ol}{\overline}
\newcommand\restr[2]{\ensuremath{\left.#1\right|_{#2}}}
\numberwithin{equation}{section}
\begin{document}
\title[Nearly invariant subspaces and weighted dual truncated Toepl]{Nearly invariant subspaces and weighted dual truncated Toeplitz operators}
\author{Sudip Ranjan Bhuia}
\date{}
\maketitle
\noindent\emph{Source and licence.} Nearly invariant subspaces and weighted dual truncated Toeplitz operators. Version arXiv:2608.03334v1 (CC BY 4.0). This material is distributed under the Creative Commons Attribution 4.0 International licence, http://creativecommons.org/licenses/by/4.0/, as recorded in the arXiv OAI licence metadata for this exact version and shown on the abstract page https://arxiv.org/abs/2608.03334v1. Full rights notice: licenses/arxiv-2608.03334-CC-BY-4.0.txt.\par\smallskip
\noindent\textit{Editorial scope.} Counted unit: the original \texttt{theorem} environment \texttt{thm:rigidity} at source lines 871--909 together with its complete proof at lines 911--1045; the delivered document keeps source lines 324--1046 so that every referenced label and every citation resolves inside it. Native editable source is the paper's own concatenated TeX; the AMS wrapper is editorial formatting only. No publisher class, upstream script or unrelated artwork is redistributed.
\begin{lemma}\label{lem:proj}
Let $\clm=h\clk_u$ be a nearly $S^*$-invariant subspace of $H^2$, where $h$ is the extremal function and multiplication by $h$ is an isometry from $\clk_u$ onto $\clm$. Let
\[
V:\clk_u\to \clm,\qquad Vk=hk ,
\]
be this unitary map. Then the orthogonal projection of $L^2(\T)$ onto $\clm$ is
\[
P_\clm=VV^*.
\]
Equivalently, for every $f\in L^2(\T)$,
\[
P_\clm f=h\,P_u(\bar h f),
\]
where $P_u(\bar h f)$ is understood through the reproducing-kernel pairing
\[
\big(P_u(\bar h f)\big)(\lambda)
=
\langle \bar h f,k_\lambda^u\rangle
=
\langle f,hk_\lambda^u\rangle_{L^2},
\qquad \lambda\in\D.
\]
Here
\[
\langle \bar h f,k_\lambda^u\rangle
:=
\int_\T \bar h f\,\overline{k_\lambda^u}\,dm
\]
is the natural $L^1$--$L^\infty$ pairing. Consequently,
\[
P_{\clm^\perp}f=f-h\,P_u(\bar h f), \qquad f\in L^2(\T).
\]
In formal notation, one may write
\[
P_\clm=M_hP_uM_{\bar h}, \qquad  P_{\clm^\perp}=I-M_hP_uM_{\bar h},
\]
provided this is interpreted in the above sense, and not as a composition of bounded multiplication operators on all of $L^2(\T)$ unless $h\in H^\infty$.
\end{lemma}

\begin{proof}
Define
\[
V:\clk_u\to L^2(\T),\qquad Vk=hk .
\]
By Hitt's theorem, multiplication by $h$ is an isometry from $\clk_u$ onto $\clm=h\clk_u$. Hence $V$ is a unitary operator from $\clk_u$ onto $\clm$, viewed as a closed subspace of $L^2(\T)$. Therefore, by the standard Hilbert space formula for the orthogonal projection onto the range of an isometry,
\[
P_\clm=VV^*.
\]

We now identify $V^*$. Fix $f\in L^2(\T)$. For every $k\in\clk_u$,
\[
\langle k,V^*f\rangle_{\clk_u}
=
\langle Vk,f\rangle_{L^2}
=
\langle hk,f\rangle_{L^2}.
\]
Thus
\[
\left|\langle hk,f\rangle\right| \leq \|hk\|_2\|f\|_2 = \|k\|_2\|f\|_2.
\]
Hence the map
\[
k\mapsto \langle hk,f\rangle
\]
is a bounded linear functional on $\clk_u$. Therefore there exists a unique element of $\clk_u$, namely $V^*f$, such that
\[
\langle k,V^*f\rangle
=
\langle hk,f\rangle
\qquad (k\in\clk_u).
\]

For $\lambda\in\D$, take $k=k_\lambda^u$. Then
\[
(V^*f)(\lambda)
= \langle V^*f,k_\lambda^u\rangle = \langle f,hk_\lambda^u\rangle.
\]
Since $h\in H^2$ and $f\in L^2(\T)$, we have $\bar h f\in L^1(\T)$. Hence the last expression may be written as
\[
\left(P_u(\bar h f)\right)(\lambda)
:= \int_\T \bar h f\,\overline{k_\lambda^u}\,dm = \langle f,hk_\lambda^u\rangle_{L^2}.
\]
Thus, in the reproducing-kernel sense,
\[
V^*f=P_u(\bar h f).
\]
Consequently,
\[
P_\clm f = VV^*f = h\,P_u(\bar h f),
\qquad f\in L^2(\T).
\]
Finally,
\[
P_{\clm^\perp}=I-P_\clm,
\]
and therefore,
\[
P_{\clm^\perp}f
=
f-h\,P_u(\bar h f),
\qquad f\in L^2(\T).
\]
This completes the proof.
\end{proof}

Throughout the paper, we regard
\[
V:\clk_u\to L^2(\T),\qquad Vk=hk,
\]
as an isometry with
\[
\operatorname{ran}V=\clm.
\]
Equivalently, $V$ is a unitary from $\clk_u$ onto $\clm$. Its adjoint is
\[
V^*:L^2(\T)\to\clk_u,
\]
and
\[
V^*V=I_{\clk_u}, \qquad VV^*=P_\clm.
\]



By extending co-analytic functions in $H_{-}^2$ via their conjugate-analytic extensions into $\D$, the orthogonal complement $\clm^\perp = \cln \oplus H_{-}^2$ forms a Reproducing Kernel Hilbert Space. The projection onto the analytic defect space $\cln$ is $P_\cln = P_{+} - P_\clm$.

\begin{lemma}\label{lem:RKHS}
Let $\clm=h\clk_u$ be a nearly $S^*$-invariant subspace and put $\cln=H^2\ominus\clm$. Then $\clm^\perp=\cln\oplus H_{-}^2$.
If
\[
f_{-}(\zeta)
= \sum_{n\geq1}a_{-n}\bar\zeta^{\,n},\qquad \zeta\in\T,
\]
define its conjugate-analytic extension to $\D$ by
\[
f_{-}^\#(\lambda)
:=\sum_{n\geq1}a_{-n}\bar\lambda^{\,n},\qquad \lambda\in\D.
\]
For $f=f_{+}+f_{-}\in\clm^\perp$, where
$f_{+}\in\cln$ and $f_{-}\in H_{-}^2$, define
\[
f(\lambda)
:=f_{+}(\lambda)+f_{-}^\#(\lambda),\qquad \lambda\in\D.
\]
Under this identification, $\clm^\perp$ is a scalar-valued harmonic reproducing kernel Hilbert space on $\D$.

Its reproducing kernel is
\[
K_\lambda^{\clm^\perp}
=
R_\lambda^\cln+k_\lambda^-,
\]
where
\[
R_\lambda^\cln(z)
=
\frac{1}{1-\bar\lambda z}
-
\overline{h(\lambda)}\,h(z)
\frac{1-\overline{u(\lambda)}u(z)}{1-\bar\lambda z},
\]
and
\[
k_\lambda^-(z)
= \sum_{n=1}^{\infty}\lambda^n\bar z^{\,n}
= \frac{\lambda\bar z}{1-\lambda\bar z}.
\]
Thus
\[
K_\lambda^{\clm^\perp}(z)
= \left(\frac{1}{1-\bar\lambda z}-\overline{h(\lambda)}\,h(z)
\frac{1-\overline{u(\lambda)}u(z)}{1-\bar\lambda z}\right)
+\frac{\lambda\bar z}{1-\lambda\bar z}.
\]
\end{lemma}

\begin{proof}
Since $\clm\subset H^2$, we have
\[
\clm^\perp
=(H^2\ominus\clm)\oplus H_{-}^2
=\cln\oplus H_{-}^2.
\]

For
\[
f_{-}(\zeta)=\sum_{n\geq1}a_{-n}\bar\zeta^{\,n}\in H_{-}^2,
\]
the series
\[
f_{-}^\#(\lambda)
=
\sum_{n\geq1}a_{-n}\bar\lambda^{\,n}
\]
converges absolutely for every $\lambda\in\D$. Indeed,
\[
|f_{-}^\#(\lambda)|
\leq
\left(\sum_{n\geq1}|a_{-n}|^2\right)^{1/2}
\left(\sum_{n\geq1}|\lambda|^{2n}\right)^{1/2}
=
\|f_{-}\|_2
\frac{|\lambda|}{\sqrt{1-|\lambda|^2}}.
\]
Hence the evaluation functionals are bounded. Moreover, the resulting
harmonic realization is injective by uniqueness of the analytic and
conjugate-analytic expansions.

The kernel for the analytic defect space $\cln$ is obtained by projecting
the Szeg\H{o} kernel $k_\lambda$ onto $\cln$. Since
\[
P_\cln=P_+-P_\clm,
\]
we have
\[
R_\lambda^\cln
=
P_\cln k_\lambda
=
k_\lambda-P_\clm k_\lambda.
\]
Using the projection formula
\[
P_\clm f=hP_u(\bar h f),
\]
where the expression is understood in the reproducing-kernel sense, we
obtain
\[
P_\clm k_\lambda
=
hP_u(\bar h k_\lambda).
\]

For every $g\in\clk_u$,
\[
\langle g,\bar h k_\lambda\rangle
=
\langle hg,k_\lambda\rangle
=
h(\lambda)g(\lambda).
\]
On the other hand,
\[
\left\langle
g,\overline{h(\lambda)}k_\lambda^u
\right\rangle
=
h(\lambda)\langle g,k_\lambda^u\rangle
=
h(\lambda)g(\lambda).
\]
Consequently,
\[
P_u(\bar h k_\lambda)
=
\overline{h(\lambda)}k_\lambda^u.
\]
Therefore
\[
R_\lambda^\cln(z)
=
k_\lambda(z)
-
\overline{h(\lambda)}h(z)k_\lambda^u(z),
\]
and hence
\[
R_\lambda^\cln(z)
=
\frac{1}{1-\bar\lambda z}
-
\overline{h(\lambda)}h(z)
\frac{1-\overline{u(\lambda)}u(z)}
     {1-\bar\lambda z}.
\]

Next, consider the co-analytic part. Since the inner product is linear
in the first variable, define
\[
k_\lambda^-(z)
=
\sum_{n=1}^{\infty}\lambda^n\bar z^{\,n}
=
\frac{\lambda\bar z}{1-\lambda\bar z}.
\]
Then
\[
\begin{aligned}
\langle f_{-},k_\lambda^-\rangle
&=
\left\langle
\sum_{n\geq1}a_{-n}\bar z^{\,n},
\sum_{n\geq1}\lambda^n\bar z^{\,n}
\right\rangle  \\
&=
\sum_{n\geq1}a_{-n}\bar\lambda^{\,n}
=
f_{-}^\#(\lambda).
\end{aligned}
\]

Finally, if
\[
f=f_{+}+f_{-}\in\clm^\perp,
\qquad
f_{+}\in\cln,\quad f_{-}\in H_{-}^2,
\]
then
\[
\begin{aligned}
\left\langle f,K_\lambda^{\clm^\perp}\right\rangle
&=
\langle f_{+},R_\lambda^\cln\rangle
+
\langle f_{-},k_\lambda^-\rangle \\
&=
f_{+}(\lambda)+f_{-}^\#(\lambda) \\
&=
f(\lambda).
\end{aligned}
\]
Hence
\[
K_\lambda^{\clm^\perp}
=
R_\lambda^\cln+k_\lambda^-
\]
is the reproducing kernel.
\end{proof}







\section{Weighted dual truncated Toeplitz operators}

Replacing the model-space complement $\clk_u^\perp$ by the orthogonal complement $\clm^\perp$ of a nearly $S^*$-invariant subspace introduces the extremal multiplier $h$ into the compression structure.

\begin{definition}\label{def:weighted_DTTO}
Let $\clm=h\clk_u$ be a nearly $S^*$-invariant subspace of $H^2$. For $\vp\in L^\infty(\T)$, the \emph{weighted dual truncated Toeplitz operator} with symbol $\vp$ is the compression of $M_\vp$ to $\clm^\perp$:
\[
D_\vp^{\clm}=\restr{P_{\clm^\perp}M_\vp}{\clm^\perp}.
\]
Equivalently,
\[
D_\vp^{\clm}f
=
P_{\clm^\perp}(\vp f),
\qquad f\in\clm^\perp .
\]
\end{definition}

\begin{remark}
Let
\[
V:\clk_u\to\clm,\qquad Vk=hk,
\]
be the canonical unitary induced by Hitt's theorem. Then
\[
P_\clm=VV^*,
\qquad
P_{\clm^\perp}=I-VV^*.
\]
Hence, for $f\in\clm^\perp$,
\[
D_\vp^{\clm}f
=
\vp f-VV^*(\vp f).
\]
Using the projection formula of Hartmann--Ross, this may be written as
\[
D_\vp^{\clm}f
=
\vp f-hP_u(\bar h\vp f).
\]
Here $\bar h\vp f\in L^1(\T)$, and $P_u(\bar h\vp f)$ is understood through the reproducing-kernel pairing
\[
\big(P_u(\bar h\vp f)\big)(\lambda)
=
\langle \bar h\vp f,k_\lambda^u\rangle
=
\langle \vp f,hk_\lambda^u\rangle,
\qquad \lambda\in\D,
\]
where the pairing is understood in the sense of
Lemma~\ref{lem:proj}. Thus the notation $hP_u(\bar h\vp f)$ represents the concrete realization of the operator $VV^*(\vp f)$.
\end{remark}

\begin{remark}
If, in addition, $h\in H^\infty$, then
$\bar h\vp f\in L^2(\T)$ for every $f\in\clm^\perp$. Using
\[
P_u=P_+-M_uP_+M_{\bar u},
\]
we obtain
\[
D_\vp^\clm f
= \vp f - hP_+(\bar h\vp f) + huP_+(\bar u\,\bar h\vp f).
\]
Since $h\in H^\infty$, all three terms belong to $L^2(\T)$.

For a general extremal multiplier $h\in H^2$, this decomposition is not
used, because the two terms involving $P_+$ need not separately belong to $L^2(\T)$. The universally valid formula is
\[
D_\vp^\clm f=\vp f-VV^*(\vp f).
\]
\end{remark}

Finally, for every $\vp\in L^\infty(\T)$,
\[
(D_\vp^{\clm})^*=D_{\bar\vp}^{\clm}.
\]
Indeed, for $f,g\in\clm^\perp$,
\[
\langle D_\vp^{\clm}f,g\rangle
=
\langle P_{\clm^\perp}(\vp f),g\rangle
=
\langle \vp f,g\rangle
=
\langle f,\bar\vp g\rangle
=
\langle f,P_{\clm^\perp}(\bar\vp g)\rangle
=
\langle f,D_{\bar\vp}^{\clm}g\rangle .
\]

\subsection{Boundedness and compactness}

Since $\clm\subset H^2$, we have
\[
\clm^\perp=(H^2\ominus \clm)\oplus H^2_{-}.
\]
In particular, $H_{-}^2\subset \clm^\perp$. This allows us to recover the classical dual Toeplitz compression as a corner of $D_\vp^\clm$.

\begin{theorem}\label{thm:boundedness}
For $\vp\in L^\infty(\T)$, the weighted dual truncated Toeplitz operator $D_\vp^\clm$ is bounded on $\clm^\perp$, and
\[
\|D_\vp^\clm\|=\|\vp\|_\infty .
\]
\end{theorem}

\begin{proof}
Since $P_{\clm^\perp}$ is an orthogonal projection, for every
$f\in\clm^\perp$ we have
\[
\|D_\vp^\clm f\|
=
\|P_{\clm^\perp}(\vp f)\|
\leq
\|\vp f\|
\leq
\|\vp\|_\infty\|f\|.
\]
Hence
\[
\|D_\vp^\clm\|\leq \|\vp\|_\infty.
\]

For the reverse inequality, note that $H_{-}^2\subset\clm^\perp$. Hence
\[
P_{\clm^\perp}P_{-}=P_{-},
\qquad
P_{-}P_{\clm^\perp}=P_{-}.
\]
Therefore, as an operator on $H_{-}^2$,
\[
P_{-}D_\vp^\clm P_{-}
=
P_{-}P_{\clm^\perp}M_\vp P_{-}
=
P_{-}M_\vp P_{-}.
\]
The operator $P_{-}M_\vp P_{-}$ on $H_{-}^2$ is the classical dual Toeplitz operator $S_\vp$. It is unitarily equivalent to a Hardy-space Toeplitz operator $T_{\widetilde{\vp}}$, where
\[
\widetilde{\vp}(z)=\vp(\bar{z}).
\]
Consequently,
\[
\|P_{-}D_\vp^\clm P_{-}\|
=
\|S_\vp\|
=
\|T_{\widetilde{\vp}}\|
=
\|\widetilde{\vp}\|_\infty
=
\|\vp\|_\infty.
\]
Since
\[
\|P_{-}D_\vp^\clm P_{-}\|\leq \|D_\vp^\clm\|,
\]
we obtain
\[
\|\vp\|_\infty\leq \|D_\vp^\clm\|.
\]
Combining the two inequalities gives
\[
\|D_\vp^\clm\|=\|\vp\|_\infty.
\]
\end{proof}

\begin{theorem}\label{thm:compactness}
For $\vp\in L^\infty(\T)$, the operator $D_\vp^\clm$ is compact on
$\clm^\perp$ if and only if $\vp=0$ almost everywhere on $\T$.
\end{theorem}

\begin{proof}
If $\vp=0$ almost everywhere, then $D_\vp^\clm=0$, and hence
$D_\vp^\clm$ is compact.

Conversely, suppose that $D_\vp^\clm$ is compact on $\clm^\perp$. Since $H_{-}^2\subset\clm^\perp$, the compression of $D_\vp^\clm$ to $H_{-}^2$ is compact. But, as above,
\[
P_{-}D_\vp^\clm P_{-}
=
P_{-}M_\vp P_{-}.
\]
Thus the classical dual Toeplitz operator
\[
S_\vp=\restr{P_{-}M_\vp}{H_{-}^2}
\]
is compact on $H_{-}^2$.

The dual Toeplitz operator $S_\vp$ is unitarily equivalent to the Hardy
Toeplitz operator $T_{\widetilde{\vp}}$, where
\[
\widetilde{\vp}(z)=\vp(\bar{z}).
\]
Therefore $T_{\widetilde{\vp}}$ is compact on $H^2$. Since a Toeplitz operator on $H^2$ with bounded symbol is compact if and only if its symbol is zero almost everywhere, we get
\[
\widetilde{\vp}=0 \quad \text{a.e. on } \T.
\]
Hence
\[
\vp=0 \quad \text{a.e. on } \T.
\]
This completes the proof.
\end{proof}




\subsection{Unitary equivalence and rigidity}

\begin{theorem}\label{thm:rigidity}
Let $\clm=h\clk_u$ be a nonzero proper nearly $S^*$-invariant subspace
of $H^2$, where $h$ is the extremal function. For
$\vp\in L^\infty(\T)$, let
\[
D_\vp
=
\restr{(I-P_u)M_\vp}{\clk_u^\perp}
\]
and
\[
D_\vp^\clm
=
\restr{P_{\clm^\perp}M_\vp}{\clm^\perp}.
\]
Then the following statements are equivalent:
\begin{enumerate}
\item The pointwise multiplication map $f\mapsto hf$ is well defined on
$\clk_u^\perp$ and defines a unitary operator
\[
W_h:\clk_u^\perp\longrightarrow\clm^\perp,
\qquad
W_hf=hf.
\]
\item The function $h$ is inner.
\end{enumerate}
When these conditions hold,
\[
W_h=\restr{M_h}{\clk_u^\perp}
\]
and
\[
D_\vp^\clm W_h=W_hD_\vp,
\qquad
\vp\in L^\infty(\T).
\]
Thus $W_h$ implements a unitary equivalence between the weighted and
classical families of dual truncated Toeplitz operators.
\end{theorem}

\begin{proof}
Suppose first that $h$ is inner. Then $|h|=1$ almost everywhere on
$\T$, so the multiplication operator
\[
M_h:L^2(\T)\longrightarrow L^2(\T)
\]
is unitary. Since
\[
M_h\clk_u=h\clk_u=\clm,
\]
unitarity gives
\[
M_h(\clk_u^\perp)=\clm^\perp.
\]
Consequently,
\[
W_h:=\restr{M_h}{\clk_u^\perp}:
\clk_u^\perp\longrightarrow\clm^\perp
\]
is unitary.

We next prove the intertwining relation. Since $M_h$ is unitary and
$M_h\clk_u=\clm$, we have
\[
P_\clm=M_hP_uM_{\bar h}.
\]
Let $f\in\clk_u^\perp$. Then
\begin{align*}
D_\vp^\clm W_hf
&=
P_{\clm^\perp}(\vp hf)\\
&=
(I-P_\clm)(\vp hf)\\
&=
\vp hf-M_hP_uM_{\bar h}(\vp hf)\\
&=
\vp hf-hP_u(|h|^2\vp f)\\
&=
\vp hf-hP_u(\vp f)\\
&=
h\bigl(\vp f-P_u(\vp f)\bigr)\\
&=
W_hD_\vp f.
\end{align*}
Therefore,
\[
D_\vp^\clm W_h=W_hD_\vp
\qquad
(\vp\in L^\infty(\T)).
\]

Conversely, suppose that the pointwise multiplication map
\[
W_h:\clk_u^\perp\longrightarrow\clm^\perp,
\qquad
W_hf=hf,
\]
is unitary. Then
\[
\|hf\|_2=\|f\|_2,
\qquad
f\in\clk_u^\perp.
\]
Hence
\[
\int_\T (|h|^2-1)|f|^2\,dm=0,
\qquad
f\in\clk_u^\perp.
\]

Since
\[
H_{-}^2\subseteq\clk_u^\perp,
\]
the preceding identity holds for every $f\in H_{-}^2$. Let $p$ be an analytic polynomial and take
\[
f(z)=\bar z\,p(\bar z).
\]
Then $f\in H_{-}^2$ and, since $|z|=1$ on $\T$,
\[
|f(z)|^2=|p(\bar z)|^2.
\]
Therefore,
\[
\int_\T (|h|^2-1)|p(\bar z)|^2\,dm=0
\]
for every analytic polynomial $p$.

Applying the polarization identity to the sesquilinear form
\[
\mathfrak q(p,q)
:=
\int_\T
(|h|^2-1)p(\bar z)\overline{q(\bar z)}\,dm,
\]
we obtain
\[
\int_\T
(|h|^2-1)p(\bar z)\overline{q(\bar z)}\,dm
=
0
\]
for all analytic polynomials $p$ and $q$. The linear span of the
functions
\[
p(\bar z)\overline{q(\bar z)}
\]
contains all trigonometric polynomials. It follows that
\[
\int_\T (|h|^2-1)\tau\,dm=0
\]
for every trigonometric polynomial $\tau$.

Since trigonometric polynomials are uniformly dense in $C(\T)$ and
\[
|h|^2-1\in L^1(\T),
\]
we conclude that
\[
\int_\T (|h|^2-1)\tau\,dm=0,
\qquad
\tau\in C(\T).
\]
Therefore,
\[
|h|^2-1=0
\quad\text{almost everywhere on }\T.
\]
Thus
\[
|h|=1
\quad\text{almost everywhere on }\T.
\]
Since $h\in H^2$ has unimodular boundary values, $h$ is inner.
\end{proof}


\end{document}
